\section{Methodology}
\label{sec:methodology}

\subsection{Japanese LMMD Construction and Validation}
\label{subsec:lmmd_construction_validation}
We construct a Japanese adaptation of the Loughran--McDonald Master Dictionary (LMMD) by translating English lexicon entries and validating the resulting Japanese terms for domain-appropriate sense and usage in EDINET filings. The workflow combines multi-engine translation (Azure, DeepL, GPT), diagnostic checks (POS alignment, corpus frequency, semantic similarity), and targeted manual review for high-impact sentiment terms.

% Previous motivation subsection
The Loughran--McDonald Master Dictionary (LMMD) is widely used in financial text analysis but is available only in English. 
Because this study examines sentiment in Japanese corporate disclosures, a direct application of LMMD is not possible. 
A cross-lingual mapping is required to transfer LMMD sentiment categories to Japanese.

Financial terminology is highly domain-specific. 
Words such as \emph{charge}, \emph{default}, \emph{bond}, \emph{provision}, or \emph{liability} exhibit substantial polysemy, and naïve machine translation often selects the incorrect sense. 
This creates a material risk that translation errors propagate into sentiment scores and bias empirical estimates. 
Accordingly, constructing a Japanese version of LMMD requires a dedicated translation, validation, and correction workflow.

% Previous Initial Machine Translation Approach subsection
As a starting point, each LMMD word was translated into Japanese using the Azure Cognitive Services Translator API. 
This produced a complete mapping from English source words to Japanese target words. 
A preliminary inspection of the translated dictionary, however, revealed several systematic deficiencies that would compromise any downstream sentiment analysis if left unaddressed.

% Previous Systematic Translation Errors subsection
Machine translation errors fell into four broad categories: 
(1) wrong sense selection, 
(2) stylistic mismatch, 
(3) transliteration or overly literal translation, and 
(4) inconsistent morphological mapping. 
These issues are illustrated in Table~\ref{tab:mt_errors}.

% Previous Wrong Sense Selection subsection
A common failure mode involved selecting the wrong semantic sense for polysemous financial terms. 
For example, the English term ``default'' was frequently translated as \jp{デフォルト}, which in Japanese typically refers to a ``preset option'' in software interfaces, rather than the correct legal and financial sense of ``credit default'' (\jp{債務不履行}). 
Similarly, ``return'' was rendered as \jp{帰る} (“to go home”), a literal verb corresponding to the everyday English meaning of “to return,” rather than the financial concept of investment return (\jp{収益率}). 
The term ``provision,'' which in accounting refers to \jp{引当金}, was translated as \jp{支給} (“payment” or “benefits”), a human-resources concept entirely unrelated to balance sheet provisions.

% Previous Stylistic Mismatch subsection
Translation engines also produced stylistically inappropriate outputs, using language that does not appear in formal corporate filings. 
For example, the English term ``leverage'' was translated as \jp{レバレッジ}, a colloquial borrowing commonly used in informal speech or media commentary. 
In Japanese financial reporting, however, leverage is typically expressed through more formal constructs such as \jp{財務レバレッジ} or \jp{有利子負債倍率}.

% Previous Transliteration and Literal Translation subsection
Some low-frequency terms were transliterated into katakana or mapped to literal everyday expressions that bear no domain relevance.
For example, ``writeoff'' was translated as \jp{償却}, which captures amortization but does not reflect the specific accounting meaning of writing off assets or receivables (e.g., \jp{貸倒償却} or \jp{評価損計上}). 
Literal or transliterated forms rarely appear in EDINET filings and therefore weaken the alignment between the Japanese lexicon and actual financial text usage.

% Previous Inconsistent Morphological Mapping subsection
A further source of noise arises from inconsistent singular–plural mappings. 
Because Japanese does not morphologically mark plurality, singular and plural English forms should ordinarily map to the same lemma. However, machine translation occasionally produced stylistically or semantically divergent outputs. 
For example, “expense’’ was translated as \jp{費用}, a general term for costs, whereas “expenses’’ was rendered as \jp{経費}, which denotes operating or administrative expenses in a narrower accounting sense. 
Although both translations are plausible, the artificial distinction between singular and plural introduces unnecessary heterogeneity into the lexicon and complicates downstream token matching.

\begin{table}[h]
\centering
\small
\caption{Machine Translation Errors in LMMD Financial Terms}
\label{tab:mt_errors}
%%\begin{tabular}{lll}
\begin{tabular}{l p{2.4cm} p{3.4cm} p{3.7cm}}
\toprule
\textbf{LMMD Term} & \textbf{MT} & \textbf{Correct Term} & \textbf{Financial Sense} \\
\midrule
default   & \jp{デフォルト} & \jp{債務不履行}    & credit default \\[3pt]
provision & \jp{支給}       & \jp{引当金}        & provision \\[3pt]
return    & \jp{帰る}       & \jp{収益率}        & rate of return \\
          &            & \jp{投資収益}      & \\[3pt]
equity    & \jp{持分}       & \jp{株主資本}      & equity capital \\
          &            & \jp{自己資本}      & \\[3pt]
leverage  & \jp{レバレッジ} & \jp{財務レバレッジ} & financial leverage \\
          &            & \jp{負債比率}        & \\[3pt]
writeoff  & \jp{償却}       & \jp{貸倒償却}        & write-off \\
          &            & \jp{評価損計上}      & \\

\bottomrule
\end{tabular}
\end{table}

These translation errors demonstrate that a naïve machine-translated lexicon cannot be reliably applied to Japanese financial disclosures without extensive validation and correction.

% Section 4 still has too many one paragraph subsections
\subsection{Translation Validation Pipeline}
\label{subsec:translation_validation}
% Previous Implications for Sentiment Analysis subsection
The quality of the sentiment lexicon directly affects classification outcomes. 
Translation errors may cause incorrect sentiment assignments, distort weighting of positive or negative terms, and reduce predictive power in empirical return regressions. 
Therefore, a more rigorous translation pipeline is required to ensure the validity of downstream results.

% Previous Improved Translation and Validation Pipeline subsection
To address the limitations of naïve machine translation, a multi-stage translation and verification procedure was implemented.

% Previous Multi-Engine Translation subsection
In addition to Azure Translate, two higher-quality translation systems were incorporated: 
DeepL, which is known for superior Japanese translation accuracy, and GPT-based models (e.g. GPT-4o and successor GPT models), which were used to refine domain-specific meanings and resolve ambiguous cases. 
For each English term, translations from multiple engines were compared for consistency.

% Previous Translation Consistency Diagnostics subsection
To assess translation stability, we translated all LMMD entries using three engines (Azure, DeepL, GPT) and compared token-level agreement. Exact match rates between engines are only 14–25\%, indicating that naïve cross-lingual dictionary transfer is highly sensitive to the translation system. We therefore apply normalization and validation steps before using the lexicon in downstream sentiment scoring.

\begin{table}[htbp]
\centering
\small
\caption{Summary statistics for Japanese LMMD translation}
\label{tab:lmmd_translation_summary}
\begin{tabular}{lrrrrr}
\toprule
\textbf{Engine} & \textbf{Rows} & \textbf{Blank} & \textbf{Nonblank} & \textbf{Valid} & \textbf{Unique JA Tokens} \\
\midrule
AZURE & 86,554 & 66  & 86,488 & 86,488 & 50,976 \\
DEEPL & 86,554 & 1   & 86,553 & 86,553 & 61,774 \\
GPT   & 86,554 & 490 & 86,064 & 86,064 & 50,270 \\
\bottomrule
\end{tabular}

\vspace{4pt}
\begin{minipage}{0.95\linewidth}
\footnotesize
\emph{Notes:} “Blank” denotes empty Japanese translations. “Valid” excludes blanks and placeholder error tokens. In the initial GPT translation run, 60 entries were returned as placeholder \texttt{ERROR\_k} tokens; these entries were subsequently corrected manually using standard Japanese equivalents.
\end{minipage}
\end{table}

\begin{table}[htbp]
\centering
\small
\caption{Pairwise exact-string agreement across MT engines}
\label{tab:lmmd_pairwise_agreement}
\begin{tabular}{lrrr}
\toprule
\textbf{Pair} & \textbf{Rows Compared} & \textbf{Exact Matches} & \textbf{Exact Match Rate} \\
\midrule
AZURE vs GPT   & 85,998 & 21,545 & 25.1\% \\
AZURE vs DEEPL & 86,487 & 18,735 & 21.7\% \\
DEEPL vs GPT   & 86,064 & 12,379 & 14.4\% \\
\bottomrule
\end{tabular}

\vspace{4pt}
\begin{minipage}{0.95\linewidth}
\footnotesize
\emph{Notes:} “Exact match” requires identical Japanese strings. Lower agreement is expected because multiple Japanese renderings can be acceptable (synonyms, formality, katakana vs.\ kanji, spacing/orthography), motivating normalization before coverage and scoring.
\end{minipage}
\end{table}

% Previous Part-of-Speech and Sense Alignment subsubsection
Part-of-speech (POS) tags from spaCy (English) and SudachiPy with UniDic (Japanese) were compared. 
A mismatch between expected POS (e.g., noun vs.\ verb) triggered a quality warning indicating probable misinterpretation.

% Previous Corpus Frequency Validation subsection
Each Japanese translation was evaluated for frequency in EDINET filings and other financial corpora. 
Correct financial translations typically appear frequently in such documents, whereas colloquial expressions or unnatural katakana forms rarely appear. 
Low-frequency translations were flagged as candidates for correction.

% PreviousSemantic Similarity Scoring subsection
Multilingual sentence embeddings (e.g., LaBSE, multilingual MiniLM, or OpenAI embeddings) were used to compute cosine similarity between the English LMMD word and its Japanese translation. 
Low similarity scores indicate potential semantic drift and were used as an additional filter.

% Previous Model-Assisted and Expert Review subsection
High-impact sentiment words (e.g., those in the negative or uncertainty categories) were manually reviewed and corrected with the assistance of GPT-based models using financial-context prompts. 
This produced a validated Japanese financial sentiment dictionary aligned with LMMD categories.

% Previous Outcome subsection
The resulting translation and validation framework produced a corrected, domain-appropriate Japanese version of LMMD. 
This lexicon enables reliable sentiment classification in Japanese corporate filings and strengthens the empirical analysis connecting textual sentiment to financial outcomes. 
Moreover, the methodology is reproducible and may be applied to sentiment lexicon construction in other languages.

\subsection{Event Study Design}
\label{subsec:eventstudy}

To evaluate whether disclosure tone is reflected in stock prices, we implement a standard short-window event study following \citet{mackinlay1997event}.

Abnormal returns are computed using a firm-specific market model. For each firm-event pair, we estimate the following regression over a pre-event estimation window:

\begin{equation}
R_{i,t} = \alpha_{i,e} + \beta_{i,e} R_{m,t} + \epsilon_{i,t},
\end{equation}

where $R_{i,t}$ denotes the daily return of firm $i$ on trading day $t$, and $R_{m,t}$ denotes the corresponding market return. The parameters $(\alpha_{i,e}, \beta_{i,e})$ are estimated via ordinary least squares using returns over trading days $[-120, -20]$ relative to the event date.

This event-specific estimation allows market model parameters to vary over time and avoids contamination from event-period returns.

Abnormal returns are then defined as:
\begin{equation}
AR_{i,t} = R_{i,t} - (\hat{\alpha}_{i,e} + \hat{\beta}_{i,e} R_{m,t}),
\end{equation}

and cumulative abnormal returns (CAR) over event window $[\tau_1, \tau_2]$ are given by:
\begin{equation}
CAR_{i,[\tau_1,\tau_2]} = \sum_{t=\tau_1}^{\tau_2} AR_{i,t}.
\end{equation}

Abnormal returns and CARs are computed independently of sentiment measures and are identical across all regression specifications.

The event date is defined as the EDINET filing date aligned to the nearest trading day on or before the filing date. We examine multiple event windows, including $[0,0]$, $[0,1]$, $[-1,1]$, $[-2,2]$, and $[-3,3]$. We take $\mathrm{CAR}[-1,1]$ as the primary event window because it is short enough to remain closely tied to the filing event while allowing for modest pre-filing information leakage, filing-time frictions, and delayed processing of detailed MD\&A language. We report $\mathrm{CAR}[0,0]$ and $\mathrm{CAR}[0,1]$ as immediacy checks, and $\mathrm{CAR}[-2,2]$ and $\mathrm{CAR}[-3,3]$ as wider-window robustness tests.

\subsection{Regression Specification and Control Variables}
\label{subsec:regression_specification}

To account for standard cross-sectional determinants of stock returns, we include firm-level control variables in extended regression specifications. Specifically, we control for firm size and return volatility.

Firm size is measured as the logarithm of market capitalization, computed using filing-date equity prices and shares outstanding obtained from EDINET XBRL disclosures. Return volatility is measured as the logarithm of the standard deviation of daily stock returns over a 60-day pre-event window ending 11 trading days prior to the filing date. Sentiment variables are standardized within the regression sample to facilitate comparison of coefficient magnitudes across sentiment measures.

To test whether disclosure sentiment is associated with market reactions, we estimate cross-sectional regressions of the form:
\begin{equation}
CAR_{i,[\tau_1,\tau_2]} =
\alpha
+ \beta Sentiment_i
+ \gamma' X_i
+ \delta_y
+ \lambda_j
+ \varepsilon_i .
\end{equation}

where $CAR_{i,[\tau_1,\tau_2]}$ denotes the cumulative abnormal return for filing $i$ over event window $[\tau_1,\tau_2]$, $Sentiment_i$ denotes the standardized sentiment measure, and $X_i$ is a vector of firm-level controls. In baseline specifications, the vector $X_i$ is omitted; in controlled specifications, $X_i$ includes log market capitalization and pre-event return volatility. The terms $\delta_y$ and $\lambda_j$ denote year and industry fixed effects, respectively. Standard errors are clustered at the firm level.

For specifications comparing GPT and LMMD sentiment directly, we estimate:
\begin{equation}
CAR_{i,[\tau_1,\tau_2]} =
\alpha
+ \beta_1 \mathrm{GPT}_i
+ \beta_2 \mathrm{LMMD}_i
+ \gamma' X_i
+ \delta_y
+ \lambda_j
+ \varepsilon_i .
\end{equation}

where $\mathrm{GPT}_i$ and $\mathrm{LMMD}_i$ denote the standardized GPT-based and Japanese LMMD-based sentiment measures for filing $i$, respectively.